 % 本文件由 scripts/assemble.py 自动装配，勿手改（改 parts/ 后重跑）
% 前言

\chapter*{前言}
\markboth{前言}{前言}

凝聚态物质（即固体和液体）的量子理论一直由两大主题主导。第一个主题是能带理论（band theory）和微扰论（perturbation theory），它大致植根于Landau的费米液体理论（Fermi liquid theory）。第二个主题是Landau的对称性破缺理论（symmetry-breaking theory）和重整化群理论（renormalization group theory）。凝聚态理论是一门非常成功的理论，它使我们能够理解几乎所有形态物质的性质。第一个主题的一项胜利是半导体理论，它为电子器件奠定了理论基础，而正是这些电子器件使近年来的技术进步成为可能。第二个主题同样重要：它使我们能够理解物质的各种状态以及它们之间的相变（phase transition），它正是液晶显示器、磁记录等等背后的理论基础。

凝聚态理论一直如此成功，以至于人们开始生出一种圆满自足之感，开始觉得自己看到的正是凝聚态理论终结的开端。然而，本书试图呈现的是另一幅图景。本书主张，我们所看到的只不过是开端的终结，在我们前方还有一整个新世界等待着我们去探索。

通往新世界的一瞥来自分数量子霍尔效应（fractional quantum Hall effect）的发现(Tsui et al., 1982)。另一瞥则来自高$T_c$超导体（high-$T_c$ superconductors）的发现(Bednorz and Mueller, 1986)。这两种现象都完全超出了上述两大主题的范畴。在过去二十年中，分数量子霍尔效应和高$T_c$超导电性方面迅速而激动人心的发展，催生了许多新思想和新概念。我们正在见证凝聚态系统多体理论（many-body theory）中一个新主题的兴起。这是凝聚态物理激动人心的时代，而这一新范式甚至可能影响我们对自然界基本问题的理解。

正是基于这一背景，我写下了本书。\footnote{1996年我开始撰写本书时，曾计划论及量子多体理论（quantum many-body theory）中一些新近而激动人心的发展。当时还不清楚这些新进展是否会成为凝聚态理论的一个新主题。而此刻，经过最近的一些进展，我本人相信一个新主题正在凝聚态理论中浮现。不过，这一理论仍处于其发展的早期阶段；我们是否真的得到了一个新主题，只有时间才能告诉我们。}本书前半部分讨论这两个旧主题，我们将其称为传统凝聚态理论（traditional condensed matter theory）。\footnote{有人也许会把第一个主题称作传统凝聚态理论，而把第二个主题称作现代凝聚态理论。}本书的第二部分则让读者得以一窥正在兴起的新主题，我们将其称为现代凝聚态理论（modern condensed matter theory）。第二部分所覆盖的材料非常新，其中一些是几个月前才问世的新结果，这一理论至今仍在迅速发展之中。

读完本书之后，我希望读者得到的不是一种圆满之感，而是一种空无之感。一百年来，凝聚态理论给予我们如此之多，而我们对自然界的丰富性仍然知之甚少。然而，我希望读者不会因此沮丧，而是为我们理解的不完整而感到兴奋。这意味着凝聚态理论那有趣而激动人心的时代还在我们前方，而非已在我们身后。我还希望读者能由此获得一种信心：没有回答不了的问题，也没有理解不了的奥秘。尽管还有许多奥秘有待理解，但我们毕竟已经理解了许多起初看似不可能理解的奥秘。我们已经理解了一些基本问题，而这些问题在起初显得太过基本，以至于似乎连答案都不会有。人类大脑的想象力同样是无边无际的。\footnote{我不知道最终谁会是“赢家”：是自然界的丰富，还是人类想象力的无边无际。}

本书源于我1996至2002年间在MIT讲授量子多体物理课程的经历。本书面向的是对现代理论物理感兴趣的研究生。第一部分（第2--5章）涵盖传统的多体物理，其中包括路径积分（path integral）、线性响应（linear response）、摩擦的量子理论（quantum theory of friction）、相互作用玻色子/费米子的平均场理论（mean-field theory）、对称性破缺与长程序（long-range order）、重整化群、正交性灾变（orthogonality catastrophe）、费米液体理论，以及非线性$\sigma$模型（nonlinear $\sigma$-model）。第二部分（第6--10章）涵盖现代多体物理中的诸多课题，其中包括分数量子霍尔理论、分数统计（fractional statistics）、流代数与玻色化（bosonization）、量子规范理论（quantum gauge theory）、拓扑序/量子序（topological/quantum order）、弦网凝聚（string-net condensation）、演生的规范玻色子/费米子（emergent gauge bosons/fermions）、量子自旋液体（spin liquid）的平均场理论，以及二维或三维的精确可解模型（exactly soluble models）。

本书所用的大多数方法都基于量子场论和路径积分。低能有效理论（low-energy effective theory）在我们的许多讨论中扮演着核心角色。即使在第一部分，我也尽量采用较为现代的方法来处理一些老问题，并尽量强调传统凝聚态物理中一些较为现代的课题。第二部分覆盖的是非常新的工作，其中约有一半来自最近几年完成的研究。第二部分的一些内容改编自我的研究论文或综述论文（而一些研究论文又是从本书的部分内容改编而来的）。

本书的写作方式着眼于强调物理图像，强调思想和观念的发展。我并不追求以干净紧凑的数学形式来呈现材料；计算和结果的表述方式旨在揭示其背后的物理图像。对于丑陋的假设，我不是把它扫到地毯下面，而是设法把它暴露出来。我还通过暴露（而非掩盖）一些常用方法所得到的错误结果，来强调这些方法的局限。

本书不打算涵盖许多不同的系统和许多不同的现象，而只讨论少数几个简单的系统。透过这几个简单的系统，我们讨论凝聚态理论中范围广泛的各种物理思想、概念和方法。书中以较小字号排印的文字是评注或更为深入的课题，初读时可以跳过。

本书的另一个特点是，我倾向于质疑并揭示多体物理、乃至更广泛的理论物理中的一些基本观念和图像，例如“什么是费米子？”“什么是规范玻色子？”这样的问题，相变与对称性破缺的观念，以及“序是否总是由一个序参量（order parameter）来描述？”等等。在这里，我们不把任何事情视为理所当然。我希望这些讨论能鼓励读者越过包裹着诸多物理思想的那些漂亮数学表述，去认识某些物理概念的丑陋与任意性。

随着数学形式体系变得越来越漂亮，人们也越发容易被形式体系所困，沦为形式体系的“奴隶”。当我们把一切都看作粒子的集合时，我们曾是牛顿定律的“奴隶”。量子理论被发现之后\footnote{经典粒子的概念在量子理论中会失效。相关讨论见2.2节。}，我们又成了量子场论的“奴隶”。眼下，我们想用量子场论解释一切，而我们的教育也不鼓励我们越过量子场论去看世界。

然而，要在物理学中取得革命性的进展，我们就不能让想象力被形式体系困住，不能允许形式体系来划定我们想象的边界。数学形式体系不过是一种工具、一门语言，它让我们得以描述想象、交流想象。有时候，当你有了一个新想法或新念头，你也许会发现自己什么都说不出来：无论说什么都是错的，因为用以描述这个新想法或新念头的确当数学或确当语言还没有被发明出来。的确，真正新颖的物理思想通常需要新的数学形式体系来描述它。这让我想起一个部落的故事。这个部落只有四个计数用的词：一、二、三和很多很多。设想部落中有一个人想到了两个苹果加两个苹果、三个苹果加三个苹果的问题，他将很难向其他部落成员解释他的理论。当你有了一个真正全新的想法时，你就应该有同样的感受。虽然本书题为《多体量子场论》（\emph{Quantum Field Theory of Many-Body Systems}），但我希望读者读完本书后会看到：量子场论并不是一切，自然界的丰富性是不会被量子场论束缚的。

我要感谢Margaret O'Meara为本书许多章节所作的校对。我还要感谢Anthony Zee、Michael Levin、Bas Overbosch、Ying Ran、Tiago Ribeiro和Fei-Lin Wang提出的意见与建议。最后但同样重要的是，我要感谢文字编辑Dr.~Julie Harris为编辑和润色本书所付出的心血。

\begin{flushleft}
\emph{Lexington, MA} \hfill Xiao-Gang Wen\\[0.5em]
\emph{October, 2003}
\end{flushleft}
